\documentclass[11pt,a4paper]{article}

\usepackage[margin=2.3cm]{geometry}
\usepackage[T1]{fontenc}
\usepackage{lmodern}
\usepackage{microtype}
\usepackage{amsmath,amssymb}
\usepackage{graphicx}
\usepackage{booktabs}
\usepackage{tabularx}
\usepackage{multirow}
\usepackage{siunitx}
\usepackage{xcolor}
\usepackage{caption}
\usepackage{subcaption}
\usepackage{enumitem}
\usepackage[section]{placeins}
\usepackage[colorlinks=true,linkcolor=blue!50!black,citecolor=blue!50!black,urlcolor=blue!50!black]{hyperref}

\graphicspath{{figures/}}
\sisetup{detect-all}
\captionsetup{font=small,labelfont=bf}
\setlist{nosep,leftmargin=*}

\newcommand{\bstar}{\ensuremath{b^\ast}}
\newcommand{\seff}{\ensuremath{\sigma_{\mathrm{eff}}}}
\newcommand{\code}[1]{\texttt{#1}}
\newcommand{\good}[1]{\textcolor{green!45!black}{#1}}
\newcommand{\bad}[1]{\textcolor{red!65!black}{#1}}

\title{Joint inflow and SGS-bias estimation in uDALES:\\
       twin tests, Round~1}
\author{pyurbanair SGS discrepancy study}
\date{October 2026}

\begin{document}
\maketitle

\begin{abstract}
\noindent
We test whether a three-coefficient strain/rotation correction to the Vreman
sub-grid-scale (SGS) viscosity in uDALES can be estimated together with the
inflow angle and speed from sparse velocity sensors. Both the truth and the
assimilation model are uDALES. The study uses parameter-only ES-MDA with 32
members over three \SI{180}{s} windows on the Xie \& Castro staggered-cube
array. Twenty-eight assimilation runs were made, plus a pilot.

\begin{itemize}
\item \textbf{Without inlet turbulence} (a near-deterministic twin), the
  coefficients are recovered on average: the T3 seed means are within 0.03 of
  the truth. This holds even though the pre-registered per-observation
  signal-to-noise gate predicted otherwise.
\item In that setting, estimating the correction lowers the held-out sensor
  and field errors by 40--90\,\% against inflow-only estimation, which is a
  clear added value.
\item The posteriors are 3--4$\times$ overconfident.
\item \textbf{With realistic inlet turbulence}, every ensemble collapses onto
  turbulence realisation noise within the first window, and the correction
  adds nothing.
\item The cause is the error model, not the correction. The observation error
  omits a realisation noise about $2.5\times$ larger than it, while 324
  observations face 32 members without localisation.
\item We also report a diagnostics bug: the posterior SGS-health values were
  copies of the prior's. It is now fixed.
\end{itemize}
\end{abstract}

\section{Questions and test design}

The correction multiplies the Vreman eddy viscosity by
\begin{equation}
  m = \exp\!\Bigl(L \tanh\frac{b_0 + b_1\,\phi(z) + b_2\,q}{L}\Bigr),
  \qquad L = \log 1.5,
  \label{eq:multiplier}
\end{equation}
where $\phi\in[0,1]$ is a sine-squared height band over
$z/H\in[0.5,2]$ (5--\SI{20}{m} for a canopy height $H=\SI{10}{m}$) and
$q\in[-1,1]$ compares the squared rotation and strain norms. The multiplier
stays within $[0.67,1.5]$, and a cell is counted as saturated when the tanh
argument exceeds about 0.74.

Round~1 asks three separate questions:
\begin{enumerate}
\item \textbf{Execution:} do the runs finish without solver failures or
  saturation?
\item \textbf{Recovery:} are the true coefficients
  $\bstar=(0.3,-0.3,0.2)$ and the true inflow (\ang{10},
  \SI{6.0}{m/s}) recovered within the posterior spread?
\item \textbf{Benefit:} does estimating $b$ improve the predictions at
  held-out sensors and the next window's forecast, compared with estimating
  the inflow alone?
\end{enumerate}

\begin{table}[!htbp]
\centering\small
\caption{Assimilation tests. Each runs in two stages: A without inlet
turbulence, B with the synthetic digital-filter inlet. L2 is the main
layout. Seeds 2 and 3 vary the assimilation seed, the prior seed and the
truth inlet seed together.}
\label{tab:tests}
\begin{tabularx}{\linewidth}{@{}llllX@{}}
\toprule
Test & Truth $b$ & Prior SGS & Estimated & Question \\
\midrule
T1 & \bstar & $\mathcal N(0,0.2^2)$ & $b$ (inflow pinned) & Is SGS identifiable at all? \\
T3 & \bstar & $\mathcal N(0,0.2^2)$ & inflow + $b$ & The main test \\
T2 & \bstar & none ($b=0$) & inflow & Uncorrected baseline for T3 \\
T0 & 0 & $\mathcal N(0,0.2^2)$ & inflow + $b$ & Does the fit invent a correction? \\
\bottomrule
\end{tabularx}
\end{table}

\paragraph{Fixed setup.}
\begin{itemize}
\item \textbf{Case:} \code{xie\_and\_castro}, 30$\times$40$\times$16 cells
  over 60$\times$80$\times$\SI{32}{m}, with cubes 2.8--\SI{17.2}{m} tall.
\item \textbf{Closure:} Vreman with $c_v=0.24$ for both truth and model.
\item \textbf{Windows:} three windows of \SI{180}{s} after a \SI{30}{s}
  spinup, with \SI{1}{Hz} output.
\item \textbf{Observations:} $u$ and $v$ averaged over \SI{20}{s}, with
  instrument error \SI{0.25}{m/s} and representation error \SI{0.1}{m/s}.
  One aggregated observation then has $\seff=\SI{0.060}{m/s}$.
\item \textbf{Smoother:} ES-MDA with 4 steps and no localisation, and
  \code{failure.policy=raise}, so a diverging member stops the run.
\item \textbf{Sensors} (Figure~\ref{fig:geometry}):
  \begin{itemize}
  \item L1: six street-level sensors at $z=\SI{2}{m}$.
  \item L2: the same six $(x,y)$ positions at $z=2$, 8 and \SI{14}{m}.
  \item VAL: four held-out validation sensors, one of them above the canopy.
  \end{itemize}
  Every sensor stencil node is a fluid cell.
\item \textbf{Configuration:} everything was set on the command line from a
  fresh clone of \code{main} (\code{06249b0}); no config was edited.
\end{itemize}

\begin{figure}[!htbp]
\centering
\includegraphics[width=\linewidth]{geometry_sensors.pdf}
\caption{Domain and sensors. (a)~Top view of the building footprints, shaded
by height (voxelised STL), with the assimilation sensors and the held-out validation sensors.
The arrow gives the true inflow direction (\ang{10}). (b)~Side view of the
building heights, with the SGS height band (5--\SI{20}{m}) where $\phi>0$ and
the sensor heights.}
\label{fig:geometry}
\end{figure}

\section{Pilot: stability and signal-to-noise}

\paragraph{Stability (P1).}
We ran single-member forward runs over three windows. Every coefficient value
tested, $b_0\in\{\pm0.3,\pm0.6\}$ and $b_1,b_2\in\{\pm0.6\}$, ran stably at
$c_v=0.24$:
\begin{itemize}
\item the minimum $\Delta t$ was \SI{0.068}{s} against \SI{0.076}{s} for
  $b=0$;
\item no cell saturated;
\item $\pm0.6$ already reaches multipliers of 0.69 and 1.44, close to the
  cap.
\end{itemize}
So $c_v=0.24$ and $\sigma_b=0.2$ were kept: $\pm3\sigma_b$ lies inside the
tested range.

\paragraph{Signal-to-noise (P2).}
Figure~\ref{fig:pilot} compares pairs of runs through
$D=\mathrm{RMS}(\Delta\text{\SI{20}{s} means})/\seff$.

The plan's gate needed $D_{\mathrm{sgs}}\ge2$ and
$D_{\mathrm{sgs}}\ge2D_{\mathrm{seed}}$, and it \bad{failed}:
\begin{itemize}
\item $D_{\mathrm{sgs}}$ was only 0.5--0.7 for \bstar.
\item The largest single coefficient, $b_0=\pm0.6$, reaches only
  $D\le1.6$, so a larger \bstar\ cannot pass under the cap.
\item With inlet turbulence on, the realisation noise is
  $D_{\mathrm{seed}}\approx3$--4, about \SI{0.2}{m/s} per \SI{20}{s} mean.
\item With it off, the realisation noise almost vanishes
  ($D_{\mathrm{seed}}\approx0.01$), but $D_{\mathrm{sgs}}$ stays at about
  0.6.
\end{itemize}

The gate is per observation, whereas one L2 window assimilates
$9\times18\times2=324$ observations jointly
($D\sqrt N\approx10$). We therefore ran the full matrix as a clearly labelled
\emph{post-gate exploratory} set, with \bstar\ unchanged.

\begin{figure}[!htbp]
\centering
\includegraphics[width=\linewidth]{pilot_snr.pdf}
\caption{Pilot signal-to-noise. $D$ is the RMS difference of \SI{20}{s}
sensor means between two runs, divided by the effective observation error
$\seff=\SI{0.060}{m/s}$, pooled over three windows. The pairs are the SGS
signal ($b=0$ against \bstar), the realisation noise (inlet seed or initial
perturbation changed) and the inflow signal (\ang{+5}, \SI{+0.5}{m/s}).
The dashed line is the gate, $D=2$.}
\label{fig:pilot}
\end{figure}

\section{Results}

\subsection{Execution}

All 24 parameter-only smoother runs finished, in 85--\SI{112}{min} each with
two side by side at 8 workers. The two state-updating methods did not:
\begin{itemize}
\item \textbf{Inlet-off joint filter: failed.} One second into its second
  cycle, just after the first joint state and parameter analysis, five
  members' $\Delta t$ fell below \SI{e-4}{s}. The diffusion number was
  pinned at 0.255 while the divergence stayed at machine precision, so the
  analysed state made the eddy viscosity blow up.
\item \textbf{Inlet-off hybrid: ran away}
  (Section~\ref{sec:methods}).
\end{itemize}

Saturation could not be judged. The posterior state files carried the
prior's SGS diagnostics (Section~\ref{sec:bug}), so the reported posterior
saturation of 14--41\,\% belongs to the prior ensemble. The posterior means
lie well inside the unsaturated range: their largest tanh argument is about
0.5, against 0.74.

\subsection{Recovery}

\begin{figure}[!htbp]
\centering
\includegraphics[width=\linewidth]{recovery_windows.pdf}
\caption{Coefficient posteriors per window, as mean $\pm$ std, for T3-L2
seeds 1--3 and T1-L2 seed 1. Lines mark the truth \bstar; the band is the
prior $\pm1\sigma$.
Top, stage~A (inlet off): the estimates approach the truth with very small
error bars, and the seed-to-seed scatter is larger than those bars.
Bottom, stage~B (inlet on): the ensembles collapse to a point in window~0,
at seed-dependent wrong values.}
\label{fig:recovery}
\end{figure}

\begin{figure}[!htbp]
\centering
\includegraphics[width=0.55\linewidth]{inflow_recovery.pdf}
\caption{Inflow posteriors per window for T3 (with SGS estimated) and T2
(inflow only), seeds 1--3; the axes are zoomed, so the wide prior bars are clipped. Lines mark the truth (\ang{10}, \SI{6.0}{m/s}).}
\label{fig:inflow}
\end{figure}

\begin{table}[!htbp]
\centering\small
\caption{Final-window posteriors on L2: the mean over seeds 1--3 of the
posterior means, with the seed-to-seed standard deviation and the mean
posterior standard deviation in parentheses.}
\label{tab:recovery}
\resizebox{\linewidth}{!}{%
\begin{tabular}{@{}llccccc@{}}
\toprule
Stage & Test & $b_0$ (0.3) & $b_1$ ($-0.3$) & $b_2$ (0.2) & angle (\ang{10}) & speed (6.0) \\
\midrule
A & T1 & 0.301 (0.052, 0.018) & $-0.276$ (0.121, 0.034) & 0.256 (0.024, 0.023) & pinned & pinned \\
A & T3 & 0.288 (0.054, 0.018) & $-0.316$ (0.136, 0.031) & 0.193 (0.051, 0.022) & 10.06 (0.15, 0.06) & 6.06 (0.01, 0.00) \\
B & T1 & 0.094 (0.137, 0.000) & $-0.094$ (0.120, 0.000) & 0.006 (0.064, 0.000) & pinned & pinned \\
B & T3 & 0.094 (0.025, 0.000) & 0.029 (0.175, 0.000) & 0.037 (0.198, 0.000) & 12.50 (1.87, 0.00) & 5.96 (0.06, 0.00) \\
\bottomrule
\end{tabular}}
\end{table}

\paragraph{Stage A: recovered on average, overconfident.}
Table~\ref{tab:recovery} and Figure~\ref{fig:recovery} give the numbers.
\begin{itemize}
\item \textbf{Accuracy.} The T3 seed means are within 0.03 of \bstar, and
  window~0 alone is close. For example, A-T1-L2 seed~1 reaches
  $(0.298,-0.339,0.162)$ after window~0.
\item \textbf{Drift.} In windows~1--2 the estimates drift while the spread
  shrinks: $b_2$ rises in every run.
\item \textbf{Overconfidence.} The posterior std is 3--4$\times$ smaller
  than the seed scatter, and the validation spread/skill ratio is only
  0.1--0.3. By the pre-registered rule
  ($|\text{mean}-\text{truth}|\le2\,\text{std}$ and
  $\text{std}\le0.5\,\sigma_b$):
  \begin{itemize}
  \item T1 passes only on L1, where the posterior is wider;
  \item it fails on L2 in all three seeds;
  \item every T3 run has at least one ``wrong'' coefficient, although the
    misses are only 0.1--0.6 prior std.
  \end{itemize}
\item \textbf{No invented correction.} The T0 control stays within 0.03 of
  $b=0$.
\item \textbf{Combinations, not single values.} The posterior correlations,
  $\rho(b_0,b_1)\approx-0.6$ and $\rho(b_0,b_2)\approx+0.5$, show that the
  sensors constrain combinations of the coefficients.
\item \textbf{Inflow.} It is collapsed and slightly biased: the speed is
  $6.06\pm0.003$\,m/s (Figure~\ref{fig:inflow}).
\end{itemize}

\paragraph{Stage B: collapse.}
\begin{itemize}
\item Every ensemble collapses to a std of about 0.001 in window~0.
\item The estimates scatter by seed with wrong signs.
\item T0 invents $b_1=+0.22$ and $b_2=-0.18$.
\item In two of three seeds the angle sticks at \ang{13.4}--\ang{13.7}.
\end{itemize}

\FloatBarrier
\subsection{Why stage B collapses}

Figure~\ref{fig:collapse} shows the mechanism. In stage~A the predicted
observations really depend on the parameters: across the ES-MDA steps, the
member spread and the ensemble-mean misfit shrink together, from 0.030 to
\SI{0.003}{m/s}. In stage~B the members' predicted-observation spread stays
at about \SI{0.125}{m/s} at every step. That spread is turbulence
realisation noise, since each member has its own inlet seed.

Two things then drive the collapse:
\begin{itemize}
\item With 324 observations, 32 members and no localisation, the 31
  parameter anomalies are always an exact linear function of the
  observation anomalies. The update therefore treats noise correlations as
  information.
\item The observation error ($\seff=\SI{0.06}{m/s}$) leaves out the
  realisation noise (0.13--\SI{0.16}{m/s} between the ensemble mean and the
  clean truth). Nothing restrains the update, and the parameter spread falls
  200-fold in one window.
\end{itemize}
The Desroziers ratio of 2.3--3.1, and the underfit flag on every stage-B
run, both say the observation error should be about three times larger.

\begin{figure}[!htbp]
\centering
\includegraphics[width=\linewidth]{collapse_mechanism.pdf}
\caption{Collapse mechanism, window~0. (a)~Median member spread of the
predicted observations, and the RMS error of the ensemble mean against the
noise-free truth, per ES-MDA step; the line marks $\seff$. (b)~Ensemble std
of $b_0$, normalised by the prior std, per window. Stage~A shrinks in step
with a genuinely informative misfit. Stage~B collapses even though its
observation spread never shrinks.}
\label{fig:collapse}
\end{figure}

\FloatBarrier
\subsection{Benefit: T3 against T2}

Figure~\ref{fig:value} and Table~\ref{tab:value} compare the held-out
scores. ``Posterior'' is the mean over all three windows; ``forecast'' is
the window-1/2 prior, a forecast from the previous posterior made before
seeing the window's data. A row counts only if $|T3-T2|$ exceeds both the
replica floor and the seed std, with the same sign in all three seeds. The
replica floor is the truth rerun with only its seed changed.

\begin{figure}[!htbp]
\centering
\includegraphics[width=\linewidth]{value_t3_vs_t2.pdf}
\caption{Relative change of held-out scores when the SGS coefficients are
estimated, $(T3-T2)/T2$ (negative is better), averaged over seeds 1--3, for
the posterior and the window-1/2 forecast. Filled markers are significant by
the rule in the text; hollow markers are within noise. The dashed line marks
the plan's $-10\,\%$ threshold.}
\label{fig:value}
\end{figure}

\begin{table}[!htbp]
\centering\small
\caption{Selected held-out scores, as means over seeds 1--3. CRPS is for
window means and variances of the velocity magnitude at the validation
sensors; W2 is the member-median Wasserstein distance; field and canopy
scores are RMSE against the truth.}
\label{tab:value}
\begin{tabular}{@{}lllllll@{}}
\toprule
Score & Stage & T2 & T3 & $\Delta$ & Replica & Verdict \\
\midrule
CRPS mean, posterior      & A & 0.0327 & 0.0109 & $-67\,\%$ & \num{6e-5} & \good{T3 better} \\
CRPS mean, forecast       & A & 0.0441 & 0.0235 & $-47\,\%$ & \num{8e-5} & \good{T3 better} \\
CRPS variance, posterior  & A & \num{6.0e-5} & \num{7.5e-6} & $-88\,\%$ & \num{2e-6} & \good{T3 better} \\
CRPS variance, forecast   & A & \num{1.47e-5} & \num{1.52e-5} & $+3\,\%$ & \num{2e-8} & T2 better (small) \\
W2 magnitude, posterior   & A & 0.0476 & 0.0186 & $-61\,\%$ & \num{1e-4} & \good{T3 better} \\
Field TKE, posterior      & A & 0.0048 & 0.0019 & $-60\,\%$ & \num{2e-4} & \good{T3 better} \\
Canopy TKE, posterior     & A & \num{4.6e-4} & \num{1.3e-4} & $-72\,\%$ & \num{1e-5} & \good{T3 better} \\
\midrule
CRPS mean, posterior      & B & 0.0846 & 0.0852 & $+1\,\%$ & 0.064 & within noise \\
CRPS mean, forecast       & B & 0.0883 & 0.0872 & $-1\,\%$ & 0.052 & within noise \\
CRPS variance, posterior  & B & 0.0131 & 0.0126 & $-4\,\%$ & 0.020 & within noise \\
Field TKE, posterior      & B & 0.0359 & 0.0358 & $0\,\%$ & 0.054 & within noise \\
\bottomrule
\end{tabular}
\end{table}

\paragraph{Stage A: adds value.}
\begin{itemize}
\item T3 is clearly better on 10 of the 12 validation CRPS rows, and all 10
  clear the $\ge10\,\%$ threshold.
\item It is also clearly better on 27 of the 34 W2, field, canopy and
  spectra rows, and clearly worse on none.
\item The plan's T3-vs-T2 pass requires $\ge10\,\%$ for the forecast
  variances too, and these fall short ($+2$ and $+3\,\%$, tiny in absolute
  terms).
\item The T3 inflow error is no larger than T2's.
\item The correctly specified T0 run gives the best reference. It scores
  0.0108 on the posterior CRPS of window means against 0.0105 for T3 and
  0.0327 for T2. Estimating $b$ thus removes essentially the whole SGS-misfit
  penalty.
\item The remaining gap to the near-zero replica floor comes from the
  biased, collapsed inflow speed.
\end{itemize}
Figure~\ref{fig:runs-canopy} shows the difference in the canopy profiles.

\paragraph{Stage B: no added value.}
T3 and T2 differ by $-4$ to $+2\,\%$, far less than the seed spread. Both
sit within $1.3\times$ the replica floor, and T0 scores the same. With
realistic inlet turbulence, the held-out scores are dominated by realisation
noise that no parameter can reduce (Figure~\ref{fig:runs-sensors}).

\begin{figure}[!htbp]
\centering
\begin{subfigure}{0.49\linewidth}
  \includegraphics[width=\linewidth]{run_A_T2_canopy_profiles.png}
  \caption{A-T2-L2 (inflow only)}
\end{subfigure}\hfill
\begin{subfigure}{0.49\linewidth}
  \includegraphics[width=\linewidth]{run_A_T3_canopy_profiles.png}
  \caption{A-T3-L2 (inflow and SGS)}
\end{subfigure}
\caption{Canopy profiles against the truth and the replica, stage~A,
seed~1 (pipeline figures from \code{visualize\_assimilation.py}).}
\label{fig:runs-canopy}
\end{figure}

\begin{figure}[!htbp]
\centering
\begin{subfigure}{0.49\linewidth}
  \includegraphics[width=\linewidth]{run_A_T3_sensor_timeseries_validation.png}
  \caption{Stage A (inlet off)}
\end{subfigure}\hfill
\begin{subfigure}{0.49\linewidth}
  \includegraphics[width=\linewidth]{run_B_T3_sensor_timeseries_validation.png}
  \caption{Stage B (inlet on)}
\end{subfigure}
\caption{Validation-sensor time series for T3-L2, seed~1: the truth against
the posterior ensemble. With inlet turbulence on, the truth's fluctuations
are a realisation the members cannot reproduce.}
\label{fig:runs-sensors}
\end{figure}

\subsection{Layout}

On stage~A seed~1, L2 beats L1 for T3 on the sensor, field and
above-canopy spectral scores; L1 is better only on the canopy-$u$ profile.
T1-L1 recovers all three coefficients, even though $b_1$ is seen there only
indirectly. Its posterior is wider (std of $b_1$ 0.08 against 0.03), so the
strict rule is easier to pass.

\subsection{Methods: smoother, hybrid, joint filter}
\label{sec:methods}

\begin{figure}[!htbp]
\centering
\includegraphics[width=\linewidth]{methods.pdf}
\caption{T3-L2 seed~1 posterior scores by method. The dashed ticks mark the
replica floor. The inlet-off joint filter failed.}
\label{fig:methods}
\end{figure}

The parameter-only smoother is the only usable method here
(Figure~\ref{fig:methods}):
\begin{itemize}
\item \textbf{Hybrid, stage~A: runs away.} Its window-0 posterior is
  reasonable, but window~1 jumps to $b_0=0.97$ and $b_2=1.35$, and window~2
  reaches $b_0=1.13$, $b_2=1.57$ and a speed of \SI{6.50}{m/s}. These values
  are far outside the prior and deep in saturation, and the CRPS is
  $10\times$ the smoother's.
\item \textbf{Hybrid, stage~B:} collapsed, like the smoothers.
\item \textbf{Joint filter, stage~B:} all members end up identical. Its TKE
  error is $6\times$ the smoother's.
\end{itemize}
The estimated coefficients do not survive the state filter.

\section{Bug: stale SGS diagnostics}
\label{sec:bug}

\code{ParameterESMDA} stacks the per-step states with
\code{xarray.concat(..., join="override")}. The default
\code{combine\_attrs="override"} stamps the prior step's attributes on the
whole stack. With \code{save\_prior\_state=true}, \code{run\_smoother.py}
writes the posterior state from that stack, so it carries the prior members'
coefficients and native multiplier diagnostics, and \code{sgs\_health}
reported the prior twice.

The fix (branch \code{fix/posterior-sgs-diagnostics}, with a regression
test) keeps each step's own attributes. The posterior values of the existing
runs need a rerun; their prior diagnostics and all posterior parameters are
correct.

Two related problems are not yet fixed:
\begin{itemize}
\item hybrid and filter forecast states keep only their first cycle's
  attributes;
\item on-disk ensembles copy member~0's attributes onto the whole ensemble.
\end{itemize}

\section{Conclusions}

\begin{description}
\item[Execution.] The parameter-only smoother is robust. The joint filter
  and the hybrid destabilise uDALES or drive the coefficients out of range.
\item[Recovery.] The correction is identifiable from sparse velocity sensors
  when realisation noise is small, despite a per-observation SNR of about
  0.6. The pre-registered gate is too conservative, and $D\sqrt N$ is the
  better predictor. The posteriors are 3--4$\times$ overconfident, however.
\item[Benefit.] Without inlet turbulence, estimating the correction
  \good{adds value}: it gives 40--90\,\% lower held-out errors and matches a
  correctly specified model. With realistic inlet turbulence it gives
  \bad{no added value}, because the ensemble collapses onto noise.
\end{description}
With 1--3 seeds per configuration, these results are evidence rather than a
statistical claim.

\paragraph{Before Round~2 (PALM truth).}
The realistic case fails on the error model, so Round~2 should not start as
planned. Test these first:
\begin{enumerate}
\item a representation error that includes the realisation noise, about
  \SI{0.15}{m/s};
\item statistics that average the noise down, such as longer aggregation or
  window means and variances, as observations;
\item localisation or a larger ensemble, to break the rank-deficient
  parameter--observation regression;
\item a rerun of the saturation diagnostics with the fix.
\end{enumerate}

\appendix
\section{Runs}

Every run used \code{ensemble\_size=32}, \code{num\_parallel\_processes=8}
and \code{ncpu=1}. The exact commands, the per-run parameter tables and the
full score tables are in the Markdown appendices next to this report
(\code{appendix\_commands.md}, \code{appendix\_params.md},
\code{appendix\_value.md}).

\begin{table}[!htbp]
\centering\small
\caption{Assimilation runs. Wall time includes the inline truth, metrics and
figures.}
\begin{tabular}{@{}lllll@{}}
\toprule
Stage & Tests & Layout & Seeds & Status \\
\midrule
A, B & T1, T3, T2 & L2 & 1, 2, 3 & all done, 85--\SI{112}{min} (two at about \SI{200}{min} under outside load) \\
A, B & T0 & L2 & 1 & done \\
A, B & T1, T3 & L1 & 1 & done \\
A, B & T3 hybrid & L2 & 1 & done, about \SI{215}{min} \\
B & T3 joint filter & L2 & 1 & done, \SI{165}{min} \\
A & T3 joint filter & L2 & 1 & \bad{failed} (diffusion-limited $\Delta t$ collapse) \\
\bottomrule
\end{tabular}
\end{table}

\end{document}
